% TOP_PROBLEM: {"id": 3223, "title": "Antichains in the cycle continuous order", "category_id": 3, "category": {"id": 3, "name": "graph_theory", "display_name": "Graph Theory", "description": "Problems involving graphs, networks, and their properties.", "slug": "graph-theory", "order_index": 3, "created_at": "2026-07-31T15:26:25.671Z"}, "status": "open", "problem_number": "OPG-323", "set_id": 12}
% FIRST_POSTED: 2026-10-01
% ENTRY_KIND: statement_only
% UnsolvedMath ID 3223; source catalogue code OPG-323
% !TeX program = lualatex
% Definition and sources only; this file is not an LLM solution attempt.
\documentclass[11pt,a4paper]{article}
\usepackage[margin=25mm]{geometry}
\usepackage{fontspec}
\setmainfont{lmroman10-regular.otf}[BoldFont=lmroman10-bold.otf,
 ItalicFont=lmroman10-italic.otf,BoldItalicFont=lmroman10-bolditalic.otf]
\usepackage{amsmath,amssymb,mathtools}
\usepackage{unicode-math}
\setmathfont{latinmodern-math.otf}
\AtBeginDocument{\let\setminus\smallsetminus}
\usepackage{xurl,hyperref}
\hypersetup{colorlinks=true,urlcolor=blue}
\setcounter{secnumdepth}{0}
\setlength{\parindent}{0pt}
\setlength{\parskip}{5pt}
\setlength{\emergencystretch}{3em}
\begin{document}
\section*{Antichains in the cycle continuous order}\label{3223}
{\small UnsolvedMath ID 3223 · OPG-323 · Graph Theory · OpenGarden}
\subsection{Definitions and mathematical statement}
For a finite undirected graph $G$, its binary cycle space $\mathcal Z(G)$ is the collection of edge subsets $F\subseteq E(G)$ for which every vertex has even degree in $(V(G),F)$. Addition is symmetric difference of subsets, making this a vector space over $\mathbb F_2$. Parallel edges are distinct coordinates; a loop contributes two to its vertex's degree. The empty subset is included, and a member may be a union of several cycles rather than one connected cycle.

Given finite graphs $G,H$, an edge map $f:E(G)\to E(H)$ is cycle-continuous if
\[
            f^{-1}(F)\in\mathcal Z(G)
            \qquad\text{for every }F\in\mathcal Z(H).
\]
No map of vertices, preservation of incidence, injectivity, or surjectivity is required. The inverse image is the ordinary set-theoretic inverse image of an edge subset; in particular every source edge mapping to an edge of $F$ is included.

The problem asks whether there is an infinite sequence of finite graphs $G_1,G_2,\ldots$, each with at least one edge, such that for every distinct $i,j$ no cycle-continuous map $E(G_i)\to E(G_j)$ exists. Since the condition is quantified over ordered pairs, it excludes maps in both directions. Such a sequence is an infinite antichain under cycle-continuous mapping. The graphs in the sequence are finite individually, with no uniform order bound required; one infinite graph is not an answer to this antichain question.
\subsection{Short English statement}
Is there an infinite family of finite graphs between which no cycle-continuous edge map exists in either direction?
\subsection{Sources}
\begin{itemize}
\item[S1] ULAM.AI and the UnsolvedMath Contributors, supplied problems.json, record 3223 (OPG-323), OpenGarden. Catalogue identity and source transcription; CC BY 4.0.
\url{https://huggingface.co/datasets/ulamai/UnsolvedMath}.
\item[S2] Open Problem Garden, Antichains in the cycle continuous order. Original problem and discussion.
\url{http://www.openproblemgarden.org/op/antichains_in_the_cycle_continuous_order}.
\end{itemize}
\end{document}
